% lecture09.tex - 基变换
% 第 9 章: 基变换 (Change of Basis)

\chapter{基变换}
\label{ch:changebasis}

\begin{abstract}
同一个向量在不同基下有不同的坐标 (coordinates)。在本章中, 我们将学习如何在不同坐标系之间转换, 以及这一概念与矩阵的关系。
\end{abstract}

% ============================================
\section{不同基下的坐标}
% ============================================

\begin{intuition}
向量的坐标 (coordinates) 依赖于所选的基。如果你说不同的语言, 同一个意思也会有不同的表达方式。

例如: 在标准基下为 $\twovec{3}{2}$ 的向量 $\vv$, 在另一组基下可能是 $\twovec{1}{1}$!
\end{intuition}

\begin{definition}[基下的坐标]
若 $\mathcal{B} = \{\vb_1, \vb_2\}$ 是一组基, 则向量 $\vv$ 在该基下的坐标 (coordinates) 是满足下式的系数 $c_1, c_2$:
\[
\vv = c_1 \vb_1 + c_2 \vb_2
\]
记号: $[\vv]_{\mathcal{B}} = \twovec{c_1}{c_2}$
\end{definition}

\begin{example}
基 $\mathcal{B} = \left\{ \vb_1 = \twovec{2}{1}, \vb_2 = \twovec{-1}{1} \right\}$

把向量 $\vv = \twovec{3}{2}$ 表示在这组基下。

需要求 $c_1, c_2$ 使得:
\[
c_1 \twovec{2}{1} + c_2 \twovec{-1}{1} = \twovec{3}{2}
\]
\[
\system{
2c_1 - c_2 &= 3 \\
c_1 + c_2 &= 2
}
\]
解为: $c_1 = \frac{5}{3}$, $c_2 = \frac{1}{3}$

于是 $[\vv]_{\mathcal{B}} = \twovec{5/3}{1/3}$
\end{example}

% ============================================
\section{过渡矩阵}
% ============================================

\begin{definition}[过渡矩阵]
从基 $\mathcal{B}$ 到标准基的\vocab{过渡矩阵 (transition matrix, 即基变换矩阵 change of basis matrix)}是以 $\mathcal{B}$ 的基向量为列构成的矩阵:
\[
\mP = [\vb_1 \,|\, \vb_2 \,|\, \cdots \,|\, \vb_n]
\]
\end{definition}

\begin{theorem}
\[
\vv = \mP \, [\vv]_{\mathcal{B}}
\]
即: 基 $\mathcal{B}$ 下的坐标乘以过渡矩阵 = 标准基下的向量
\end{theorem}

\begin{theorem}
\[
[\vv]_{\mathcal{B}} = \mP\inv \vv
\]
即: 要从标准基转换到基 $\mathcal{B}$, 需要用逆矩阵。
\end{theorem}

% ============================================
\section{基变换的变换视角}
% ============================================

\begin{intuition}
矩阵 $\mP$ 是一个恒等变换, 它只改变语言:
\begin{itemize}
    \item $\mP$: 把语言 $\mathcal{B}$ 翻译成标准语言
    \item $\mP\inv$: 把标准语言翻译成语言 $\mathcal{B}$
\end{itemize}
\end{intuition}

\begin{center}
\begin{tikzpicture}[scale=1.2]
    % 标准基
    \begin{scope}[shift={(0,0)}]
        \draw[grid] (-0.5,-0.5) grid (3.5,2.5);
        \draw[axis] (-0.5,0) -- (3.5,0);
        \draw[axis] (0,-0.5) -- (0,2.5);
        \draw[basis] (0,0) -- (1,0) node[below] {$\vi$};
        \draw[basis] (0,0) -- (0,1) node[left] {$\vj$};
        \draw[vector, blue!70, very thick] (0,0) -- (3,2) node[above] {$\vv$};
        \node at (1.5,-1) {标准基};
        \node at (1.5,-1.5) {$\vv = \twovec{3}{2}$};
    \end{scope}

    % 另一组基
    \begin{scope}[shift={(6,0)}]
        \draw[gray!30] (-0.5,-0.5) -- (1,0.5) -- (3,1.5) -- (1.5,1) -- cycle;
        \draw[axis] (-0.5,0) -- (3.5,0);
        \draw[axis] (0,-0.5) -- (0,2.5);
        \draw[red!70, very thick, -{Stealth}] (0,0) -- (2,1) node[below right] {$\vb_1$};
        \draw[red!70, very thick, -{Stealth}] (0,0) -- (-1,1) node[above left] {$\vb_2$};
        \draw[vector, blue!70, very thick] (0,0) -- (3,2) node[above] {$\vv$};
        \node at (1.5,-1) {基 $\mathcal{B}$};
        \node at (1.5,-1.5) {$[\vv]_{\mathcal{B}} = \twovec{?}{?}$};
    \end{scope}
\end{tikzpicture}
\end{center}

% ============================================
\section{在不同基下表示变换}
% ============================================

\begin{theorem}[新基下的变换矩阵]
若 $\mA$ 是变换 $T$ 在标准基下的矩阵, 则同一变换在基 $\mathcal{B}$ 下的矩阵为:
\[
[\mA]_{\mathcal{B}} = \mP\inv \mA \mP
\]
\end{theorem}

\begin{intuition}
这个公式包含三步:
\begin{enumerate}
    \item $\mP$: 从基 $\mathcal{B}$ 翻译到标准基
    \item $\mA$: 在标准基下实施变换
    \item $\mP\inv$: 把结果翻译回基 $\mathcal{B}$
\end{enumerate}
\end{intuition}

\begin{center}
\begin{tikzpicture}[node distance=3cm]
    \node (B1) {$[\vv]_{\mathcal{B}}$};
    \node (S1) [right of=B1] {$\vv$};
    \node (S2) [right of=S1] {$T(\vv)$};
    \node (B2) [right of=S2] {$[T(\vv)]_{\mathcal{B}}$};

    \draw[-{Stealth}, thick] (B1) -- (S1) node[midway, above] {$\mP$};
    \draw[-{Stealth}, thick] (S1) -- (S2) node[midway, above] {$\mA$};
    \draw[-{Stealth}, thick] (S2) -- (B2) node[midway, above] {$\mP\inv$};
    \draw[-{Stealth}, thick, blue] (B1) to[bend right=30] node[midway, below] {$\mP\inv\mA\mP$} (B2);
\end{tikzpicture}
\end{center}

% ============================================
\section{基变换的重要性}
% ============================================

\begin{intuition}
为什么基变换 (change of basis) 重要?

有些变换在特定基下看起来\textbf{更简单}。例如:
\begin{itemize}
    \item 在标准基下旋转很复杂
    \item 但在一条轴落在旋转轴上的基下, 它就变简单了
\end{itemize}

对一个变换最好的基? \textbf{特征基 (eigenbasis)}---下一章见!
\end{intuition}

\begin{practical}
\textbf{应用: 简化计算}

假设要计算 $\mA^{100}$。如果在合适的基下, $\mA$ 变为对角矩阵:
\[
\mP\inv \mA \mP = \mD \quad \text{(diagonal)}
\]
则:
\[
\mA^{100} = \mP \mD^{100} \mP\inv
\]
而 $\mD^{100}$ 非常容易计算!
\end{practical}

% ============================================
\section{练习}
% ============================================

\begin{exercise}
已知基 $\mathcal{B} = \left\{ \twovec{1}{1}, \twovec{1}{-1} \right\}$, 求向量 $\vv = \twovec{3}{1}$ 在该基下的坐标 (coordinates)。
\end{exercise}

\begin{exercise}
写出从 $\mathcal{B} = \left\{ \twovec{2}{1}, \twovec{1}{1} \right\}$ 到标准基的过渡矩阵 (transition matrix)。
\end{exercise}

\begin{exercise}
若变换在标准基下的矩阵为 $\mA = \twomat{2}{1}{0}{3}$, 求该变换在基 $\mathcal{B} = \left\{ \twovec{1}{0}, \twovec{1}{1} \right\}$ 下的矩阵。
\end{exercise}

\begin{exercise}[挑战]
证明 $\det(\mP\inv\mA\mP) = \det(\mA)$。
\end{exercise}

\begin{problem}
已知两组基 $\mathcal{B}_1$ 和 $\mathcal{B}_2$。如何直接计算从 $\mathcal{B}_1$ 到 $\mathcal{B}_2$ 的过渡矩阵 (不经过标准基)?
\end{problem}
